\section{Methods}

This is a secondary analysis of released docking scores and published assay tables. We ran no
new docking and made no new measurement. OpenAI Codex assisted with code review, implementation
review of sensitivity analyses, literature search and language editing. Anthropic Claude (Fable 5)
was used for an independent adversarial audit of the manuscript and its claim boundaries. The authors
executed the analyses, inspected the outputs and verified the reported values; the full disclosure is
in the Declarations. Seed contracts, tie-resolution rules, bin contracts
and the detailed null constructions are in the Supplementary Information.

\subsection{Score matrices and preprocessing}

Docking-44 is the frozen export of Nikitin et al.~\cite{nikitin2025}: 12{,}651 ligands by 44
Vina-GPU~2.0 target columns \cite{ding2023}. The carrier file also contains three appended score
columns---8pm6 (BSEP), 6itm (FXR) and 9fzj (PXR)---that are not among the 44 targets in the
published source panel. We excluded them by matching target identities to that panel before any
score statistic was calculated. The receptor structures are inherited from that
study and are not uniformly human: five of the 43 Protein Data Bank (PDB) entries carry a non-human receptor chain,
which matters wherever a human sequence baseline is compared against them, and we report a
construct-clean sensitivity for the panel where it applies. Within the selected 44 columns, we
clipped positive scores to zero, which changed nothing because none is strictly positive, and
replaced the 8{,}159 missing cells (1.47\% of
556{,}644) by their target-column mean. DOCKSTRING-58 is the DOCKSTRING v1 table
\cite{dockstring2022}, 260{,}155 molecules by 58 AutoDock Vina columns \cite{vina2010}. We kept
the 260{,}060 rows that are complete across all 58 targets, dropped the 95 rows carrying any of
the 339 missing cells, and clipped the 5{,}393 positive scores to zero. Everything ran in the
pinned release environment: Python 3.12.11 with NumPy 1.26.4, pandas 2.3.2 and SciPy 1.16.2
\cite{numpy2020,scipy2020}.

\subsection{Target maps, centring and spectral summaries}

For an $N\times P$ score matrix $X$, the raw target map is the sample Pearson correlation matrix
of its $P$ target columns. The residual surface is the two-way-centred matrix
\begin{equation}
\Gamma_{ij}=X_{ij}-\bar X_{i\cdot}-\bar X_{\cdot j}+\bar X_{\cdot\cdot},
\label{eq:centering}
\end{equation}
and the residual map is the target correlation matrix of $\Gamma$. Correlation already removes
column means, so the only operative difference between the two maps is that the residual map has
each ligand's mean across targets taken out. A cell of $\Gamma$ asks a relative question: not
how well this compound scored here, but how well it scored here compared with how it scored
across the panel.

We summarise a map by the eigenvalues $\lambda_k$ of its correlation matrix.
Participation-ratio (PR) dimension is $(\sum_k\lambda_k)^2/\sum_k\lambda_k^2$. It counts how
many directions the correlations effectively spread over: $P$ if the targets were mutually
uncorrelated, close to 1 if they all moved as one. Since $\sum_k\lambda_k=P$ and
$\sum_k\lambda_k^2=P+2\sum_{j<k}r_{jk}^2$, it is exactly
\begin{equation}
r_{\mathrm{PR}}=\frac{P}{1+(P-1)\,\overline{r_{jk}^{2}}},
\label{eq:pr}
\end{equation}
where the bar averages over the $P(P-1)/2$ off-diagonal target correlations. PR therefore adds
no information beyond the correlation spectrum. It restates the mean squared off-diagonal
correlation on a scale readers can compare with a target count, and it is not an algebraic
rank. PC1 variance fraction and entropy effective rank are secondary summaries
\cite{roy2007,mazzucato2016}. The shared-axis diagnostic orients the leading eigenvector to a
positive loading sum and takes its cosine with the unit-normalised all-ones vector, so that a
value near 1 means the leading direction is close to ``all targets rise together''. Agreement
between two maps is the Spearman correlation of their corresponding strict-upper-triangle
entries.

Mode selection takes the smallest $k$ whose leading residual eigenvalues reach half the residual
trace, read from the docking spectrum with no experimental value entering. The retained map is
the rank-$k$ reconstruction rescaled to unit diagonal, and $k$ was recomputed in five
scaffold-disjoint folds.

\subsection{Fixed-pose rescoring across scoring functions}

To separate the scoring function from the pose search, we rescored one fixed block of retained
DOCKSTRING poses: 2{,}000 ligands by 9 targets, 18{,}000 cells, each holding the top AutoDock
Vina pose. Because every scorer sees the same geometry, no difference between their outputs
can be attributed to a different selected pose. The outputs still differ in implementation,
molecular and receptor perception, training set and software environment, so what the
comparison isolates is scorer-and-implementation sensitivity at fixed geometry, not the
behaviour of independent end-to-end docking pipelines. smina 2020.12.10 supplied
the Vina and Vinardo functions \cite{smina2013,vinardo2016}, ODDT 0.7 supplied RF-Score v1, v2
and v3, NNScore 2.0 and PLECscore linear
\cite{oddt2015,rfscore2010,ballester2014rfscore,nnscore2011,plecscore2019}, and the released
DOCKSTRING Vina column is the eighth output. RF-Score and NNScore were trained on PDBbind 2016
descriptor tables of 3{,}767 complexes; PLECscore uses the coefficients ODDT ships from that
release. Scores were not clipped here, because clipping only the Vina-family columns would treat
the scorers asymmetrically.

The nine targets form a convenience block spanning GPCRs, nuclear receptors, enzymes and one
kinase, with retained poses available for every sampled ligand. They were not randomly sampled from
DOCKSTRING and are not claimed to represent its full target panel.

Four criteria were evaluated for each output: PC1 fraction at least 0.50; cosine at least 0.95 to
the uniform target direction with every loading positive; raw PR no larger than half the target
count; and a centring-induced PR rise whose conservative cluster-bootstrap interval excludes
zero. They are a descriptive checklist. The criteria and the results they judge enter the
repository in one commit, so their ordering cannot be verified from the released package.

\subsection{Transformation-matched nulls}

The residual dimension has to be read against something, because the subtraction raises it
whatever the data. Each matrix therefore has one deterministic simple-random support of 12{,}000
ligands, and every null family within a matrix uses that support and 500 repetitions. Every
simulated surface passes through the same centring and spectral steps as the observed one, so
the only thing that varies is the mechanism that generated the numbers. The additive null adds
Gaussian noise with the observed target-specific residual standard deviations to fitted grand,
ligand and target effects. The column-permutation null permutes each residual target column
independently across ligands, and the direction null replaces each ligand's residual by a random
zero-sum direction with its norm preserved. The size-conditioned null permutes each raw target
column within ten equally populated molecular-weight bins, which keeps each target's score
distribution inside a coarse size domain while destroying any correspondence between targets for
a given ligand.

\subsection{Chemical-domain controls}

Molecular weight came from RDKit on the frozen analysis SMILES. The low band satisfies
$\mathrm{MW}\le q_{0.25}$ and the high band $\mathrm{MW}\ge q_{0.75}$, with boundary ties kept in
both, and each band has its own raw and row-centred map. The matched control is 200 splits of
the same library into size-matched halves that share no whole chemical group, which asks how far
two maps drift apart when the chemistry differs but the size distribution does not. Groups are
the frozen \texttt{Butina\_clusters} column for Docking-44, 8{,}150 clusters \cite{butina1999},
and RDKit Bemis--Murcko scaffolds for DOCKSTRING \cite{murcko1996}, with every acyclic molecule a
singleton. Docking-44 used all 12{,}651 rows, DOCKSTRING a fixed 15{,}000-row complete-case
sample.

Calibration recovery used 100 outcome-blind samples without replacement at each of seven sizes
from 50 to 5{,}000 ligands, comparing each sample's residual map with the full-matrix map. Its
chemical-group version splits the library into five \texttt{GroupKFold} folds under the same
rule, drawing 25 samples of 200 and of 500 ligands from four folds and taking the reference map
from the fifth.

\subsection{Experimental resources and matching}

The kinase panels are DAVIS, 72 compounds as $\mathrm{p}K_\mathrm{d}$
\cite{davis2011,wu2025davis}; PKIS2, 645 compounds over 21 mapped percentage-inhibition columns
\cite{drewry2017}; PKIS1, 360 compounds \cite{elkins2016}; and KiRHub, 92 drugs by 20 targets,
which gives names but no structures \cite{saifudeen2026}. Secondary pharmacology comes from the
PDSP $K_\mathrm{i}$ database, 481 compounds over a 27-target export whose primary graph keeps the
148 target pairs with at least ten shared compounds across 21 targets
\cite{pdspdatabase,pdsp2000,helmuth2000}, and the
Novartis Safety Profiling Database, 12 targets and 1{,}975 InChIKeys with 91.6\% of selected rows
right-censored \cite{sutherland2023}.
PDSP combines measurements from different publications and assay conditions. Median aggregation
within a compound--target cell reduces duplicate rows but cannot remove study or batch effects.

Compound identity is the Standard InChIKey recomputed from source SMILES with the pinned RDKit
\cite{rdkit2025}, and matching an assay compound to a docking row required the full key to be
equal. Exclusion is deliberately coarser than matching, so that near-relatives leave as well:
any docking row whose 14-character connectivity block appears in an assay panel left the
calibration support first. That dropped 481 DOCKSTRING rows over 1{,}050 blocks for the kinase
panels, and 1{,}364 rows over 932 blocks for the safety panel, which shares no exact key with
DOCKSTRING at all. For PDSP, the conservative mask used every exact or censored compound retained
by the analysis and removed 128 Docking-44 rows before its map was built. The exact-only endpoint
itself shared 63 full keys and 116 connectivity blocks with Docking-44.

The kinase endpoint is defined by the assay panels alone, with no docking score entering: a
target pair is positive when its centred experimental correlation falls in the upper decile of at
least two of DAVIS, PKIS2 and PKIS1, which labels 15 of the 190 pairs on the 20 shared kinases.
Structural baselines on the same pairs are receptor-domain sequence identity from a global
BLOSUM62 alignment, and identity across the aligned 85-residue Kinase--Ligand Interaction
Fingerprints and Structures (KLIFS) ATP pocket \cite{klifs2021}.
Centring an assay matrix changes what its edges measure, so every raw-versus-residual comparison
holds the experimental map fixed and varies only the docking side.

\subsection{Statistics}

Edges of a map share nodes: two target pairs that involve the same receptor are not independent.
Inference therefore permutes whole target labels rather than individual edges. A quadratic
assignment procedure permutation test \cite{hubert1976} applies one complete relabelling of the
docking targets to the rows and columns of its map while the experimental endpoint stays fixed,
and paired residual-minus-raw contrasts reuse the same relabelling, so the two maps are always
judged against the same shuffle. Permutation and simulation probabilities use the add-one rule,
with 50{,}000 permutations on the kinase and safety-panel endpoints and 20{,}000 on the PDSP
graph. The external comparisons are post hoc. The directional kinase and safety-panel tests retain
the one-sided positive alternatives defined in their source analyses, but we treat their
probabilities as descriptive tail areas. For PDSP, the paired-QAP
probabilities reported in the paper are two-sided: each is twice the smaller of the add-one
upper- and lower-tail probabilities, capped at one;
directional values remain in the released table as diagnostics only. We report Holm and Bonferroni
adjustments within each 12-row PDSP QAP table: two contrasts, three metrics and two permutation
schemes. These within-table adjustments do not correct the wider post-hoc analytical search. QAP preserves
dependence among edges that share targets, but it does not propagate the uncertainty in estimating
each PDSP edge from a different set of compounds. Where a separate kinase contrast
was reported on both ROC-AUC and average precision, we adjusted across those two metrics by Holm. A
homology-preserving version of the same test shuffles labels only within receptor families, which asks
whether an association survives once relatedness among receptors is held fixed.

Ligand-side uncertainty comes from a cluster bootstrap that resamples chemical groups with
replacement rather than individual ligands. The groups are Butina clusters built from Morgan
fingerprints or whole Bemis--Murcko scaffolds
\cite{morgan1965,rogers2010,murcko1996}. Fixed-pose intervals are the conservative
union of the two schemes at 5{,}000 replicates each. All of these are finite-support summaries on
fixed target panels. They are not confidence intervals for a superpopulation of targets.
